\section[The q20 FAB calibration]{The failed FAB calibration for the recurrent state model}
\label{sec:bf-fab-q20-failure}

The September 28, 2026 calibration identified three distinct difficulties:
biased replay sampling on the saved buffers, inadequate evidence of accurate
AIS gradients, and a non-integrable auxiliary target for the mathematical UKF
posterior at the initial IAF checkpoint. The last finding changes the next
scientific question. Before asking how well a Markov kernel samples
$\gamma^2/q$, one must establish that this function can be normalized.
Faithfully porting the author's transitions does not establish that assumption
for a new posterior and proposal.\footnote{\raggedright
\def\UrlBreaks{\do\/\do\-}
The calibration results, source-anchored tail audit, arithmetic checks and
executed source are preserved under
\path{plans/artifacts/fab-monograph-coverage-2026-09-28/evidence/}.
The manuscript plan
\path{plans/bayesfilter-fab-monograph-coverage-2026-09-28.md}
records their provenance.\par}

The application called q20 has twenty latent coordinates and twenty hidden
coordinates, a sixty-dimensional augmented state, and thirty observations.
Only four parameters are inferred: a latent mean weight and bias, and an
observation weight and bias. Its prior is $\calN(\mu,16I_4)$, with
$\mu=(0.35,-0.08,0.65,0.05)$. The frozen initial proposal uses the canonical
three-stage IAF, two width-sixteen ELU hidden layers per stage, conditional
log-scale cap two, a standard Gaussian base, and outer affine scale four.
These are the actual study settings, not recommendations for other targets.

\subsection{Replay bias is visible without sampling the posterior}

Fix a saved map and its entire $3{,}232$-row replay buffer. Let $w_i$ denote
the normalized stored priority, $c_i=q_{\rm old}(x_i)/q_\phi(x_i)$, and
$h_i=-\nabla_\phi\log q_\phi(x_i)$. The declared capped finite-buffer
expectation is exactly computable:
\begin{equation}
 G_{\rm buf}=\sum_{i=1}^{3232}w_i\min(c_i,10)h_i.
 \label{eq:bf-fab-buffer-target}
\end{equation}
This quantity is conditional on the buffer and the cap; it is not the exact
FAB gradient. The source sampling rule first selects $128$ distinct rows,
shuffles them, and uses $32$ for the first update. If row $i$ has inclusion
probability $\rho_i$ in that subset, its expected coefficient in the first
minibatch mean is $\rho_i/128$, rather than $w_i$. In particular, that
coefficient cannot exceed $1/128$, however large the row's priority. Drawing
$32$ indices independently with replacement according to $w$ instead has
expectation \eqref{eq:bf-fab-buffer-target}.

At each of two saved maps, $256$ independent first minibatches were compared
with the exact buffer sum, without optimizer updates. The mean-gradient
relative errors were as follows; brackets give approximate conditional
bootstrap error intervals.
\begin{center}
\small
\begin{tabular}{lcc}
\toprule
Saved map & Distinct-row rule & Replacement rule \\
\midrule
Seed 1 & $0.99565\ [0.96053,1.03077]$
       & $0.01922\ [0,0.06054]$ \\
Seed 2 & $0.99265\ [0.97972,1.00558]$
       & $0.00881\ [0,0.02618]$ \\
\bottomrule
\end{tabular}
\end{center}
This supports a substantial bias of the distinct-row rule for these buffers.
The without-replacement bias is acknowledged in
\citet[Section 3.2]{midgley2023fab}; reproducing the heuristic correctly does
not ensure that its bias is small. Replacement repairs this particular
finite-buffer expectation, but not the buffer's coverage. Many stale density
corrections exceeded the cap, and refreshing all priorities with their
uncapped corrections gave an ESS near one. Full training would also have to
handle repeated indices, parameter changes and priority refresh correctly.
The experiment did not compare complete optimization trajectories.

\subsection{Finite AIS diagnostics did not qualify a training setting}

Five settings completed eight independent measurement batches of size $32$
after pilot tuning. Here $K$ counts interior bridge distributions, excluding
the endpoints, and $L$ counts leapfrog steps per HMC transition. Measurements
used frozen steps, the same initial IAF, and the same target. The map used
FP32 and the target FP64, with GPU/XLA execution, TF32 disabled, and verified
incremental GPU allocation.

For batch $b$, write $\widehat G_b=-\sum_i\bar W_{bi}
\nabla_\phi\log q_\phi(X_{bi})$, and let $\bar G$ be the mean across the
eight batches. The noise column below is the square root of the sample
variance trace of $\widehat G_b$, divided by $\|\bar G\|$. The difference
column is $\|\bar G-G_{\rm ref}\|/\|G_{\rm ref}\|$; the final column is
the root mean squared difference of individual batches from $G_{\rm ref}$,
on the same relative scale. These diagnostics remove the author's extra
factor $1/B$ in the fresh training loss so that changing batch size does not
merely rescale the comparison.
\begin{center}
\small
\begin{tabular}{lrrr}
\toprule
Setting & Relative noise & Mean difference & RMS difference \\
\midrule
Metropolis, linear $K=10$ & 1.41773 & 0.31971 & 1.41926 \\
HMC $L=5$, linear $K=10$ & 1.37817 & 0.47606 & 1.61843 \\
HMC $L=9$, linear $K=10$ & 0.37932 & 0.86541 & 1.00209 \\
HMC $L=5$, geometric $K=10$ & 0.82318 & 0.50710 & 1.05473 \\
HMC $L=5$, geometric $K=21$ & 0.66453 & 0.53804 & 0.95960 \\
\bottomrule
\end{tabular}
\end{center}
The $K=21$ setting used bridge-weighted acceptance in its pilots. The $L=9$
screen combines two original batches and six independent continuation
batches with the same frozen settings; device contention prevents a fair
runtime ranking. The finite reference used $32{,}768$ fresh independent prior
draws, reweighted by $\gamma^2/(qr)$, where $r$ is their prior density. Its
ESS was only $77.39$, and its approximate relative bootstrap gradient radius
was $0.3797$. It was therefore a noisy finite bank, not posterior truth or an
exact auxiliary expectation.

None of these measurements met the study's predeclared total relative
accuracy allowance of $0.5$. That allowance was a conservative screening
choice, not a theorem about what noise stochastic optimization can tolerate.
The observed differences do not establish a ranking of AIS settings.
In particular, the smaller measured variance for $L=9$ coexisted with a
large difference from the finite reference. Ordinary acceptance could also
conceal rejection among important particles: at the final three interior
temperatures of the linear $L=5$ screen, bridge-weighted median acceptance
was $0.025/0.098/0.014$, whereas ordinary mean acceptance was
$0.604/0.564/0.467$. Weight concentration makes these weighted estimates
noisy, but the contrast shows why acceptance alone is insufficient.

The refinement to $K=43$ completed four pilots and two measurement batches.
It was deliberately stopped during its third measurement after the tail
finding below, before its time cap expired. A generic worker status called
the interruption a timeout; the termination record identifies the mathematical
validity decision. It is not a sixth completed comparison or an observed
numerical crash. The study charged $17{,}722.765$ GPU seconds, or about
$4.923$ of its six allocated GPU hours. It trained no new map.

\subsection{Why the intended auxiliary integral diverges}

Let $\gamma^0_{\rm UKF}(x)=r(x)L^0_{\rm UKF}(x)$ denote the ordinary
real-valued UKF likelihood times the Gaussian prior. The superscript zero
specifies a principal positive-semidefinite square-root extension when a
state covariance is singular. This mathematical value target is distinguished
below from the executable's covariance margins and score-status restrictions.
The likelihood is approximate relative to the underlying nonlinear model;
the following calculation concerns this declared UKF approximation.

Consider a tube, rather than a single ray,
\begin{equation}
 x(t,u)=\mu+(u_0,u_1,t,t+u_3),\qquad
 t\ge t_0,\quad |u_0|,|u_1|,|u_3|\le10^{-6}.
 \label{eq:bf-fab-tail-tube}
\end{equation}
Its fixed width is a positive-volume witness, not a proposed parameter bound.
The latent weight and bias remain in a compact set. Each LSTM hidden output
is a sigmoid gate times a hyperbolic tangent and is therefore bounded in
absolute value by one. Consequently the latent predictions are uniformly
bounded over the tube and incoming state sigma points, even when some cell
state coordinates become large.

The augmented UKF rule has dimension $80$, fixed strictly positive diagonal
innovation covariance $Q$, and parameters $\alpha=1$, $\beta_{\rm UKF}=2$,
$\kappa=0$. Its off-center covariance weights are $1/160$ and its center
covariance weight is two, all nonnegative. An innovation-axis pair has latent
predictions $f(m)\pm\sqrt{80}\sqrt Qe_j$, because the augmented placement
covariance is block diagonal. Put $d=f(m)-\bar f$ and
$\delta_j=\sqrt{80}\sqrt Qe_j$. This pair contributes
\begin{align}
 &\frac1{160}\{(d+\delta_j)(d+\delta_j)^\top
                 +(d-\delta_j)(d-\delta_j)^\top\}\nonumber\\
 &\qquad=\frac{dd^\top}{80}+\sqrt Qe_je_j^\top\sqrt Q.
 \label{eq:bf-fab-innovation-bound}
\end{align}
Summing these pairs and the other nonnegative contributions gives
$P_{\rm pred,latent}\succeq Q$. The observation row $H(t)$ has first
coefficient $0.65+t$, with other coefficients fixed. For its scalar predictive
variance $S(t)=H(t)P_{\rm pred,latent}H(t)^\top+R$, $R>0$, there are positive
constants, uniform over the tube and the thirty dates, such that
\begin{equation}
 c_1(1+t)^2\le S(t)\le c_2(1+t)^2.
 \label{eq:bf-fab-observation-tail}
\end{equation}
The lower bound follows from $Q$; bounded latent sigma points give the upper
bound. The predictive observation mean is $O(t)$, including its growing bias.
Since the recorded observations are finite, the squared standardized
prediction error is bounded. Each Gaussian predictive density is therefore
bounded below by a positive constant times $(1+t)^{-1}$, giving
\begin{equation}
 L^0_{\rm UKF}(x(t,u))\ge C(1+t)^{-30}.
 \label{eq:bf-fab-likelihood-lower}
\end{equation}
The recursion exists at every finite $t$: nonnegative covariance weights
give a positive-semidefinite joint moment matrix, adding $R>0$ makes the
observation variance positive, and its Schur complement leaves a
positive-semidefinite filtered covariance for the next principal square root.

For the saved map, the represented neural weights were treated as exact real
numbers in an independent interval calculation of the sequential inverse.
ELU preserves a positive-slope affine asymptote and approaches $-1$ on a
negative-slope input. The capped scales approach finite constants. Resolving
every nonzero branch slope uniformly over the tube cross-section gives
\begin{align}
 T_\phi^{-1}(x(t,u))&=a(u)t+b(u)+o(1),\\
 186.3135886007891&<\|a(u)\|^2<186.3135888648922.
 \label{eq:bf-fab-inverse-tail}
\end{align}
The remainder is uniform, and $a,b$ are bounded over that cross-section.
Outward-rounded calculations at eighty and 160 digits agree; 138 independent
arithmetic enclosure checks support the calculation. Finite-ray plots alone
would not establish this bound on a region of positive volume.

At this checkpoint the twelve free scale biases are zero, so the forward
log determinant lies between $4\log4-24$ and $4\log4+24$. Combining the
Gaussian base density with \eqref{eq:bf-fab-inverse-tail} gives
\begin{equation}
 -\log q_\phi(x(t,u))\ge
 \tfrac12(186.3135886007891)t^2-C_4t-C_5.
\end{equation}
Meanwhile, twice the log prior is
\begin{equation}
 2\log r(x(t,u))=-\frac{t^2}{8}-\frac{u_3t}{8}
 -\frac{u_0^2+u_1^2+u_3^2}{16}+\text{constant}.
\end{equation}
The tube center in the original FP32 affine translation differs slightly
from the displayed decimal center. Recentering the tube on that represented
translation preserves the inverse calculation and changes only linear and
constant terms in this prior expansion. For either translation,
\begin{equation}
 \log\frac{\gamma^0_{\rm UKF}(x(t,u))^2}{q_\phi(x(t,u))}
 \ge 93.0317943003945\,t^2-C_6t-60\log(1+t)-C_7.
 \label{eq:bf-fab-tail-divergence}
\end{equation}
The quadratic coefficient is strictly positive. The map from $(t,u)$ to $x$
has absolute Jacobian one and a fixed positive cross-section volume.
Integration of this lower bound over the tube proves
\begin{equation}
 \int_{\R^4}\frac{\gamma^0_{\rm UKF}(x)^2}{q_\phi(x)}\,dx=\infty.
 \label{eq:bf-fab-q20-infinite}
\end{equation}
Thus the normalized auxiliary law and the normalized-gradient expectation
in \eqref{eq:bf-fab-gradient} do not exist for this target and checkpoint.
The posterior itself is proper: the observation-noise floor bounds the
likelihood above by a constant, and the Gaussian prior is integrable.
Posterior properness and FAB integrability are different requirements.

The executable adds small covariance margins, finite-precision repairs and
strict covariance/score checks; invalid rows are returned as NaN. The audit
does not prove that the whole unbounded tube passes those checks, or that an
explicitly restricted density $\gamma\mathbf1_{\rm valid}$ has the same
divergence. Such a restriction would be a different target whose support and
normalization have not been justified. A numerical failure cannot silently
become a declaration of zero posterior probability. The proved failure is
therefore that of the ordinary UKF value target; the global semantics of the
status-restricted numerical target remain unresolved.

This result rules out treating further finite AIS banks as convergence to the
normalized auxiliary law of that mathematical target. It does not prove that
remote tails caused every earlier finite-time training failure, that every
IAF has inadequate tails, or that FAB fails on targets satisfying its
assumptions. The next repair must first specify a well-defined sampling and
training objective. Section~\ref{sec:bf-is-multimode-recovery} develops one
that preserves the IAF while avoiding this inverse-density tail requirement.
